% !TEX root = ../main.tex

\chapter{全文总结与展望}

\section{主要结论}

本文围绕“基于执行反馈的强化学习代码生成研究”这一主题，构建了从数据整理、监督微调到 PPO 强化学习优化的完整实验流程，并在 MBPP 和 HumanEval 数据集上进行了实验验证。本文的核心思想是利用代码任务天然具备的可执行特性，将模型生成代码后的运行结果、测试通过情况和错误信息转化为奖励信号，从而使模型优化目标更接近代码生成任务的功能正确性要求。

本文首先对 MBPP 和 HumanEval 数据进行了统一处理，将原始数据转换为适用于 SFT 和 PPO 的两类格式。SFT 数据主要包含 prompt 和 response，用于监督微调；RL 数据保留 prompt、参考答案、入口函数和测试代码，用于生成、执行和 reward 计算。最终实验数据包含 MBPP train 385 条、MBPP valid 42 条和 HumanEval test 164 条。为了保证执行反馈可靠，本文对标准答案进行了 runtime 复核，确认所有样本的测试代码和入口函数均可被执行环境正确处理。

在模型训练方面，本文以 Qwen2.5-Coder-1.5B 作为基座模型，采用 LoRA 进行监督微调，得到 SFT baseline。随后，本文构建了代码生成、执行测试、reward 计算、advantage 估计和 PPO policy update 的闭环流程。该流程能够读取 RL 格式数据，使用当前策略模型生成代码，对生成结果进行清洗和入口函数对齐，在本地执行环境中运行测试，并根据执行反馈计算奖励。PPO 更新阶段则通过 clipped policy loss、value loss 和 KL penalty 对 LoRA adapter 进行进一步优化。

实验结果表明，本文实现的 PPO 方法能够在 MBPP valid 上改善 SFT 后模型的表现。在统一评测口径下，SFT baseline 在 rl\_valid 上的 Pass@1 为 0.5714，PPO best 在 reward\_v2 设置下达到 0.6667，追平 Base 模型的 Pass@1。该结果说明，在当前验证集分布下，基于执行反馈的 PPO 能够修复小样本 SFT 后模型出现的性能下降，使模型生成代码更符合测试执行要求。

同时，实验也表明，Qwen2.5-Coder-1.5B 本身已经具备较强的代码生成能力。Base 模型在 rl\_valid 上的 Pass@1 已经达到 0.6667，在 HumanEval 测试集上的 Pass@1 为 0.5000。这说明小规模 SFT 和 PPO 并不必然全面超过强代码基座模型，而更适合作为面向特定任务分布的策略调整方法。本文实验中，SFT 在 HumanEval 上相对 Base 有小幅提升，Pass@1 从 0.5000 提升到 0.5244；PPO best 在 HumanEval 上的 Pass@1 为 0.5122，略低于 SFT baseline，但 execution failures 降为 0。该结果说明 PPO 提升主要集中在验证集分布上，跨数据集泛化能力仍有限。

在 reward 设计方面，本文比较了原始测试通过率奖励、reward\_v2、reward\_v3 和 reward\_v3b。实验表明，reward\_v2 是当前最有效的奖励函数。该奖励将代码可执行、测试通过比例和全部测试通过三种信号组合起来，使 PPO 能够获得比原始 reward 更密集的执行反馈。在 reward\_v2 下，PPO best 的 valid Pass@1 达到 0.6667。相比之下，reward\_v3 和 reward\_v3b 虽然引入了 AST 相似度和函数签名匹配等结构奖励，但 Pass@1 均为 0.5952，未能超过 reward\_v2。这说明在本文小数据和单步 PPO 设置下，结构相似性奖励与最终功能正确性之间存在偏差，执行测试反馈仍是更可靠的优化信号。

在 PPO 调参方面，本文比较了 learning rate、KL 系数和 rollout 数量的影响。实验发现，单独降低 KL 系数会导致策略偏移过大，使模型表现退化；单独降低 learning rate 可以使更新更保守，但收益也随之减弱；将 rollout 数量从 42 扩大到 100 并没有自动带来性能提升。当前最佳设置为 42 条 rollout、learning rate 为 $1\times10^{-5}$、KL 系数为 0.02、batch size 为 4，并使用 reward\_v2。该结果说明，PPO 效果依赖于 reward 密度、rollout 规模和 update 强度之间的匹配。

综合来看，本文的主要结论可以概括为以下几点：第一，执行反馈可以有效补充 SFT 的 token 级训练目标，使代码生成优化更接近功能正确性；第二，在小规模数据条件下，PPO 能够在验证集分布上改善 SFT 后模型表现，但不一定稳定提升跨数据集泛化能力；第三，奖励函数设计是影响 PPO 效果的关键因素，简单、直接且与测试通过率强相关的 reward\_v2 比复杂结构奖励更有效；第四，PPO 训练需要合适的 KL 约束和更新强度，更多 rollout 或更复杂 reward 并不必然带来更好结果。

\section{本文工作的局限性}

尽管本文完成了从 SFT 到 PPO 的完整代码生成强化学习实验闭环，并取得了一定实验结果，但当前工作仍存在若干局限。

首先，本文使用的数据规模较小。SFT 和 PPO 主要基于 MBPP 数据，其中训练样本为 385 条，验证样本为 42 条。虽然这些数据已经足以验证方法流程和 reward 设计，但相对于代码大模型训练和强化学习后训练所需的数据规模仍然较小。小规模数据容易导致模型对特定题型和测试风格产生局部适配，限制其跨数据集泛化能力。HumanEval 测试结果也表明，PPO best 未能超过 SFT baseline，说明当前训练分布与测试分布之间仍存在差异。

其次，本文实现的 PPO 框架仍是简化版本。当前方法将一次完整代码生成视为单步 episode，虽然包含 policy loss、value loss、KL penalty 和单步 advantage，但尚未实现完整多步 trajectory PPO。对于语言模型生成任务而言，代码是逐 token 生成的长序列，完整建模每一步 token 的动作、状态和回报传播会更加复杂。本文的单步近似降低了实现难度，也适合有限算力条件下的实验，但其优化能力和稳定性仍弱于完整大规模 PPO 训练框架。

再次，reward 设计仍存在进一步优化空间。reward\_v2 在验证集上取得了最好结果，但它仍主要依赖当前测试用例的通过情况。如果测试用例覆盖不足，模型可能学到通过当前测试的局部策略，而不是获得真正稳健的程序语义理解。reward\_v3 和 reward\_v3b 尝试加入 AST 相似度和函数签名匹配，但实验表明简单结构奖励并未提升 Pass@1。这说明如何设计既密集又与功能正确性高度一致的 reward，仍然是代码生成强化学习中的关键问题。

此外，当前执行环境相对简化。本文采用本地 Python 执行环境完成生成代码测试，能够满足 MBPP 和 HumanEval 的基本实验需求，但还不具备完整沙箱系统的资源隔离和超时控制能力。对于更复杂的代码生成任务，如果生成代码包含死循环、文件操作、网络请求或高资源消耗行为，简单本地执行方式可能存在安全性和稳定性问题。后续若扩展到更复杂数据集，需要引入更严格的沙箱执行机制。

最后，本文主要采用 Pass@1 和 average reward 作为核心评价指标。虽然这些指标能够反映代码是否通过测试，但对代码可读性、复杂度、鲁棒性和安全性等方面的评价仍不充分。实际软件工程场景中的代码质量不仅取决于是否通过当前测试，还涉及可维护性、异常处理、性能和安全性等因素。因此，本文结果更适合说明执行反馈对函数级代码生成正确性的影响，而不能完全代表实际工程代码质量。

\section{研究展望}

针对上述局限，后续研究可以从以下几个方向继续推进。

第一，扩大训练数据规模并提高数据多样性。后续可以引入更多代码生成数据集，如扩展版 MBPP、EvalPlus、APPS 或真实开源项目中的函数级任务，使模型在更丰富的问题类型和测试分布上进行训练。更大规模和更多样化的 rollout 数据有助于缓解当前 PPO 只在验证集分布上提升的问题，并可能提升模型在 HumanEval 等测试集上的泛化能力。

第二，改进 reward 设计。本文实验表明，reward\_v2 比原始 reward 和结构 reward 更有效，但仍然依赖测试用例覆盖。后续可以考虑结合多种反馈信号，例如隐藏测试生成、错误类型分类、静态分析、代码复杂度约束、边界条件覆盖率等，构建更稳健的 reward。同时，也可以研究分阶段 reward，将语法正确性、入口函数一致性、部分测试通过和边界测试通过分别建模，降低奖励噪声。

第三，完善 PPO 训练框架。后续可以实现更完整的多步 PPO，将 token 级生成过程纳入轨迹建模，并引入更稳定的 value function 训练、GAE 估计和 KL 自适应控制。对于更大模型和更大数据集，还可以使用分布式 rollout 和批量执行环境，提高训练效率和样本利用率。

第四，增强执行环境的安全性和鲁棒性。代码生成模型可能产生不可控代码，因此执行反馈训练需要安全隔离。后续可以将本地执行器替换为具备超时控制、资源限制和文件系统隔离能力的沙箱环境，避免死循环或危险操作影响训练过程。同时，可以记录更细粒度的错误类型，为 reward 设计和案例分析提供更多信息。

第五，扩展评价指标和案例分析。除 Pass@1 外，后续可以引入 Pass@k、平均测试通过率、代码复杂度、运行时间和更严格的隐藏测试集评价。对于模型退化案例，可以进一步分析 PPO 更新前后的输出分布变化，研究哪些类型的问题更容易被执行反馈改善，哪些类型的问题容易发生边界条件退化。

总体而言，基于执行反馈的强化学习为提升代码生成模型的功能正确性提供了一条可行路径。本文在小规模数据和有限算力条件下验证了该路径的基本可行性，也揭示了 reward 设计、训练稳定性和泛化能力方面的挑战。未来随着更高质量数据、更完善执行环境和更稳定强化学习算法的发展，执行反馈有望在代码生成模型后训练中发挥更重要作用。
